%  Workshop No. 2 - Technical Architecture and Bronze/Silver Pipeline
%  Volcanic Activity Monitoring System (United States), Team 6
%
%  Build: pdflatex report.tex (twice). Screenshots go in workshop_2/figures/ with the names
%  used below; a missing image prints a framed placeholder instead of breaking the build.
\documentclass[11pt,a4paper]{article}

\usepackage[utf8]{inputenc}
\usepackage[T1]{fontenc}
\usepackage[english]{babel}
\usepackage{lmodern}
\usepackage{textcomp}
\usepackage[a4paper,margin=2.3cm]{geometry}
\usepackage{graphicx}
\usepackage{booktabs}
\usepackage{longtable}
\usepackage{array}
\usepackage{caption}
\usepackage{enumitem}
\usepackage{xcolor}
\usepackage{listings}
\usepackage{fancyhdr}
\usepackage[hidelinks]{hyperref}

\graphicspath{{figures/}{./}{../docs/}}

\hypersetup{
  pdftitle={Workshop 2 - Volcanic Activity Monitoring: Technical Architecture and Bronze/Silver Pipeline},
  pdfauthor={Team 6 - Universidad Distrital Francisco Jose de Caldas},
  pdfsubject={Data Analysis Programming 2026-II, Workshop No. 2}
}

\captionsetup{font=small,labelfont=bf}
\setlist[itemize]{itemsep=2pt,topsep=3pt}
\setlist[enumerate]{itemsep=2pt,topsep=3pt}

\pagestyle{fancy}
\fancyhf{}
\lhead{\small Workshop 2 --- Volcanic Activity Monitoring}
\rhead{\small Team 6}
\cfoot{\thepage}
\renewcommand{\headrulewidth}{0.4pt}

\newcolumntype{L}[1]{>{\raggedright\arraybackslash}p{#1}}

\lstdefinestyle{evidence}{
  basicstyle=\scriptsize\ttfamily, breaklines=true, columns=flexible,
  frame=single, framesep=4pt, rulecolor=\color{gray!50}, backgroundcolor=\color{gray!5}
}
\lstset{style=evidence}

% screenshot that may not exist yet
\newcommand{\shot}[3]{%
  \begin{figure}[htbp]\centering
  \IfFileExists{figures/#1}{\includegraphics[width=#2\linewidth]{#1}}%
    {\fbox{\parbox[c][4cm][c]{0.8\linewidth}{\centering\small Screenshot pending: \texttt{figures/\detokenize{#1}}}}}
  \caption{#3}\label{fig:#1}
  \end{figure}}

\begin{document}

\begin{titlepage}
\centering
{\large Universidad Distrital Francisco José de Caldas\\ School of Engineering --- Computer Engineering\par}
\vspace{1.2cm}
{\Large Data Analysis Programming 2026-II\par}
\vspace{1.5cm}
{\huge\bfseries Workshop No.~2\\[4pt] Technical Architecture and Bronze/Silver Pipeline\par}
\vspace{0.8cm}
{\Large Volcanic Activity Monitoring System (United States)\par}
\vspace{1.5cm}
{\large Team 6\par}
\vspace{0.4cm}
\begin{tabular}{ll}
20231020006 & Juan Sebastián Rangel Pinto\\
20231020055 & Jhonatan David Vergara Jiménez\\
20231020064 & Julián Pulido\\
\end{tabular}
\vfill
{\large Professor: Eng. Carlos Andrés Sierra, M.Sc.\par}
\vspace{0.3cm}
{\large November 2026\par}
\end{titlepage}

\tableofcontents
\newpage

\section{Executive summary}

In Workshop~1 we chose a United States scope for volcanic activity monitoring and validated our sources with real samples: Smithsonian GVP volcanoes and eruptions, the USGS earthquake feed, and the USGS Hazard Notification System (HANS) alert levels. This workshop turns that design into a running pipeline that follows the Medallion pattern:

\begin{itemize}
  \item \textbf{Stack.} Docker Compose runs Apache Airflow~2.10 (webserver on \texttt{localhost:8080}, scheduler, triggerer) with a PostgreSQL~16 metadata database. The project folders are mounted as volumes, so DAG edits are picked up live and every output stays on the host.
  \item \textbf{Bronze.} The \texttt{bronze\_ingestion} DAG (TaskFlow API, hourly) calls three APIs (HANS, the USGS feed, and GVP once a month). It also \emph{scrapes} the USGS Volcano Updates web page with BeautifulSoup. Each source goes through extract, validate and persist. The payload is stored as \texttt{<source>\_YYYYMMDD\_HHMMSS.json}, and nothing is written when the content did not change.
  \item \textbf{Silver.} The \texttt{silver\_processing} DAG starts with a \texttt{FileSensor} on \texttt{datalake\_bronze/}. It cleans, types, validates and deduplicates the new files with pandas, and writes one Parquet file per topic (\texttt{<topic>\_YYYYMMDD\_HHMMSS.parquet}). Rejected records go to a quarantine Parquet, and every row points back to its Bronze file.
  \item \textbf{Results.} The first runs turned 5 Bronze files into 165 volcanoes, 1,281 US eruptions, 54 alert notices, 229 earthquakes within 30~km of a US volcano and 9 scraped notices. 144 records went to quarantine. Running both DAGs again produced no duplicate rows. 62 unit tests run on recorded fixtures.
\end{itemize}

Every preprocessing rule is grounded in a data issue observed in the Workshop~1 samples or in these first runs (Section~\ref{sec:prep}).

\section{Architecture}

\begin{figure}[htbp]\centering
\IfFileExists{architecture_w2.png}{\includegraphics[width=\linewidth]{architecture_w2.png}}%
  {\fbox{\parbox[c][4cm][c]{0.8\linewidth}{\centering\small Render \texttt{workshop\_2/architecture.mmd} to \texttt{workshop\_2/figures/architecture\_w2.png}}}}
\caption{Pipeline architecture. Solid arrows run today; dotted arrows are the Gold layer (Workshop~3) and the dashboard (Workshop~4).}
\label{fig:arch}
\end{figure}

\begin{longtable}{L{2.6cm}L{2.6cm}L{3.6cm}L{5.6cm}}
\toprule
\textbf{Component} & \textbf{Technology} & \textbf{Input / output} & \textbf{Purpose and interactions}\\
\midrule
\endhead
Data sources & REST APIs, WFS, HTML page & --- / HTTP responses & HANS (JSON), USGS feed (GeoJSON), GVP WFS (GeoJSON), USGS Volcano Updates (HTML). No API keys; polite User-Agent from \texttt{.env}.\\
Source clients & Python, \texttt{requests}, BeautifulSoup + lxml & HTTP / bytes + parsed payload & \texttt{src/volcano\_pipeline/sources/}: timeouts, 4 retries with exponential backoff on 429/5xx, \texttt{Retry-After}. Reused from Workshop~1, plus the new scraper.\\
Airflow scheduler & Airflow 2.10, LocalExecutor & DAG files / task runs & Triggers both DAGs on schedule and runs their tasks as local subprocesses.\\
Airflow webserver & Airflow 2.10 & metadata DB / UI on 8080 & Monitoring, manual triggers, logs.\\
Airflow triggerer & Airflow 2.10 & --- & Required service for deferrable operators; kept so the stack matches the standard Airflow layout.\\
Metadata DB & PostgreSQL 16 & Airflow state & DAG runs, task instances, XComs, connections. Not used for project data in this workshop.\\
Bronze layer & JSON files & source payload / \texttt{datalake\_bronze/} & Immutable, one file per source and run, self-describing envelope (source, time, URL, SHA-256, payload).\\
Silver layer & pandas + PyArrow, Parquet & Bronze JSON / \texttt{datalake\_silver/} & Typed, validated, deduplicated tables per topic, with lineage columns; quarantine in \texttt{datalake\_silver/quarantine/}.\\
Gold layer (planned) & PySpark (local mode) & Silver folder / \texttt{datalake\_gold/} & Workshop~3: governance KPIs and analytical summaries.\\
Dashboard (planned) & Plotly Dash & Gold & Workshop~4: Volcano Activity Indicator for the functional user.\\
\bottomrule
\end{longtable}

\subsection{Technology choices}

\begin{itemize}
  \item \textbf{Airflow with the TaskFlow API} gives scheduling, retries, logs and a UI without extra code; \texttt{@dag}/\texttt{@task} keep the DAG files short. Business logic lives in \texttt{src/volcano\_pipeline}, so it is unit-tested without Airflow.
  \item \textbf{LocalExecutor instead of Celery.} The workshop text mentions a worker. On one laptop a Celery worker also needs Redis, and it adds memory pressure: Docker Desktop has 3.7~GB here. LocalExecutor runs tasks as scheduler subprocesses with the same DAG code, and moving to CeleryExecutor later only changes the compose file.
  \item \textbf{pandas over Polars} for Silver: batches are small (the largest Bronze file is the 11~MB GVP eruptions layer), the team already used pandas in Workshop~1, and the PyArrow backend gives the Parquet schema control we need.
  \item \textbf{Parquet for Silver:} columnar, compressed (the 11.6~MB GVP eruptions JSON becomes a 130~KB Parquet with the 1,281 US eruptions), typed (time zones, categories, nullable integers survive a round trip), and readable as a folder by PySpark in Workshop~3.
  \item \textbf{JSON envelope for Bronze:} the workshop asks for JSON. The envelope keeps provenance (URL, time, HTTP status, SHA-256 of the original bytes) inside the file, so a Bronze file can be re-processed on its own.
\end{itemize}

\subsection{File nomenclature}

\begin{tabular}{ll}
\toprule
Bronze & \texttt{datalake\_bronze/<source\_id>\_YYYYMMDD\_HHMMSS.json} (UTC retrieval time)\\
Silver & \texttt{datalake\_silver/<topic>\_YYYYMMDD\_HHMMSS.parquet} (UTC processing time)\\
Quarantine & \texttt{datalake\_silver/quarantine/rejected\_YYYYMMDD\_HHMMSS.parquet}\\
Silver state & \texttt{datalake\_silver/\_state/manifest.json}, \texttt{seen\_<topic>.parquet}\\
\bottomrule
\end{tabular}

\section{Data sources in the pipeline}

\begin{longtable}{L{3.2cm}L{2.3cm}L{2.0cm}L{6.8cm}}
\toprule
\textbf{Source id} & \textbf{Access} & \textbf{Schedule} & \textbf{Notes}\\
\midrule
\endhead
\texttt{usgs\_hans} & API, JSON & hourly & \texttt{volcano/getMonitoredVolcanoes}: latest notice for 69 monitored volcanoes.\\
\texttt{usgs\_feed} & API, GeoJSON & hourly & \texttt{all\_week.geojson} (1,800+ events worldwide). A 7-day window, so hourly polling loses nothing.\\
\texttt{gvp\_volcanoes} & WFS, GeoJSON & monthly & 165 US Holocene volcanoes (\texttt{CQL\_FILTER=Country='United States'}); GVP v5.4.0.\\
\texttt{gvp\_eruptions} & WFS, GeoJSON & monthly & Full layer (11,089 eruptions); US subset selected by volcano number in Silver.\\
\texttt{usgs\_volcano\_updates} & \textbf{web scraping}, HTML & hourly & \url{https://www.usgs.gov/programs/VHP/volcano-updates}: server-rendered page with one block per elevated volcano and the latest notice of each observatory (AVO, CalVO, CVO, HVO, NMI, YVO). \texttt{robots.txt} allows \texttt{/programs/}.\\
\bottomrule
\end{longtable}

\section{Preprocessing pipeline (Bronze to Silver)}\label{sec:prep}

The complete rule set, with the observation behind each rule, is in \texttt{workshop\_2/preprocessing.md}. Summary:

\subsection{Schema enforcement and data types}
Each topic has an explicit Arrow schema (\texttt{silver/schemas.py}). Writes add missing columns as nulls, drop unknown columns and cast types, so all files of a topic share one schema. Times are UTC timestamps. The sources use three encodings: epoch milliseconds in the feed, epoch seconds in HANS, and partial dates in GVP. Low-cardinality text such as alert level, colour code, magType and network becomes \texttt{category}; volcano numbers, years and VEI are nullable integers.

Two type issues appeared only with real data. First, eruption dates before 1677 overflow nanosecond timestamps (the first run failed on an eruption in 1510), so \texttt{start\_time} uses microseconds. All Parquet timestamps are written with microsecond precision anyway: Spark~3, which reads this folder in Workshop~3, rejects nanosecond Parquet timestamps (found in the final code review). Second, the scraped notices give local times such as \emph{10:48 AM AKST}, which are converted with an explicit table of US time zones.

\subsection{Null handling}
\begin{longtable}{L{3.5cm}L{4.2cm}L{6.6cm}}
\toprule
\textbf{Field} & \textbf{Observed} & \textbf{Strategy}\\
\midrule
\endhead
Eruption VEI & null in 480 of 1,281 (37\%) & keep NULL + \texttt{vei\_missing}; never imputed (a guessed explosivity would bias the indicator)\\
Eruption month/day & 0 or null & first day of the period + \texttt{time\_precision}\\
BCE eruptions & 387 & no timestamp, year kept in \texttt{start\_year}\\
\texttt{Last\_Eruption\_Year} & null (``Unknown'' in older exports) & NULL\\
Feed magnitude & nullable per USGS docs & keep row + \texttt{mag\_missing}\\
Feed geometry & required for the volcano join & quarantine \texttt{missing\_coordinates}\\
HANS \texttt{vnum} null & 2 region-wide notices & quarantine \texttt{region\_level\_notice}\\
HANS without notice & 13 volcanoes & quarantine \texttt{no\_notice\_issued}\\
HANS volcano not in GVP list & e.g. Yellowstone & keep + \texttt{coords\_missing}\\
\bottomrule
\end{longtable}

\subsection{Validation, outliers and deduplication}
Ranges: latitude $[-90,90]$, longitude normalised to $[-180,180)$, depth $[-5,800]$~km (negative depths are valid at volcanoes), magnitude $[-2,10]$. Non-earthquake events (129 quarry blasts, explosions and ice quakes this week) are quarantined. Only earthquakes within 30~km of a US volcano with $M\ge1.0$ are kept; the 1,453 others are counted, not stored.

Outliers in magnitude and depth are flagged with the IQR rule ($1.5\times$IQR) and the modified z-score ($|0.6745(x-\tilde{x})/\mathrm{MAD}|>3.5$)~\cite{iglewicz}, and \emph{never removed}: swarms and shallow events are the signal at volcanoes. First run: 3 and 2 magnitude flags, 31 and 62 depth flags.

Deduplication works at three levels. Inside a batch, the last version of each \texttt{record\_id} wins. Across runs, a seen-index keeps the latest content hash per record. Bronze files already processed are listed in a manifest. Without the seen-index, each feed event would be stored up to 168 times, because the feed returns the same 7 days every hour. The same mechanism keeps the quarantine from growing with repeated rejects.

\subsection{Text preprocessing}
The scraped notices are free text, so they go through the NLP step:
\begin{enumerate}
  \item Strip HTML: the real text contains \texttt{<p>}, \texttt{<strong>}, \texttt{<sub>} and \texttt{\&nbsp;}.
  \item Unicode NFKC normalisation, which keeps names such as Kīlauea and Pāhala.
  \item Lowercase and remove URLs and punctuation.
  \item Tokenise and remove English stop words from a fixed list.
  \item Set keyword flags by word stem (eruption, ash, lava, seismicity, deformation, gas, plume).
\end{enumerate}
The most frequent tokens were \emph{volcanic, activity, episode, eruption, fountaining, november, summit, lava}. Words such as \emph{november} justify a domain stop-word list in Workshop~3.

\section{Docker stack}

\begin{tabular}{lll}
\toprule
\textbf{Service} & \textbf{Image} & \textbf{Role}\\
\midrule
\texttt{postgres} & postgres:16 & Airflow metadata (named volume \texttt{airflow-metadata})\\
\texttt{airflow-init} & volcano-airflow:2.10.5 & migrates the DB, creates the UI user, exits\\
\texttt{airflow-webserver} & volcano-airflow:2.10.5 & UI on port 8080\\
\texttt{airflow-scheduler} & volcano-airflow:2.10.5 & schedules and executes tasks (LocalExecutor)\\
\texttt{airflow-triggerer} & volcano-airflow:2.10.5 & deferrable operators\\
\bottomrule
\end{tabular}

\medskip
The custom image extends \texttt{apache/airflow:2.10.5-python3.11} with \texttt{airflow/requirements.txt} (requests, PyYAML, python-dateutil, beautifulsoup4, lxml, pandas, pyarrow), installed with the official constraints file. The bind mounts are \texttt{airflow/dags}, \texttt{airflow/logs}, \texttt{airflow/plugins}, \texttt{src} (on \texttt{PYTHONPATH}), \texttt{config}, \texttt{datalake\_bronze}, \texttt{datalake\_silver} and \texttt{datalake\_gold}. All services share the default Compose network, and only port 8080 is published. The \texttt{FileSensor} connection \texttt{fs\_default} is defined through the environment variable \texttt{AIRFLOW\_CONN\_FS\_DEFAULT}.

\begin{lstlisting}
cp .env.example .env            # set HTTP_USER_AGENT with a team contact
docker compose up -d --build    # first build takes a few minutes
# open http://localhost:8080  (user airflow / password airflow by default)
# unpause bronze_ingestion and silver_processing in the UI
\end{lstlisting}

\section{DAG documentation}

\subsection{\texttt{bronze\_ingestion}}
\begin{itemize}
  \item \textbf{Schedule:} \texttt{@hourly}, no catch-up, one active run; 3 retries with exponential backoff.
  \item \textbf{Tasks:} \texttt{clean\_tmp} first, then three task groups, each running \texttt{extract} $\rightarrow$ \texttt{validate} $\rightarrow$ \texttt{persist}:
  \texttt{api\_source} (mapped over HANS and the feed), \texttt{gvp\_source} (mapped over the two GVP layers; does nothing if the layer was already checked this month) and \texttt{scrape\_source} (web scraping). The groups are mapped per source, so a failed source does not block the others.
  \item \textbf{Validation:} expected structure per source; for the scraper, at least one volcano block must be found, so a page redesign fails the task (\texttt{AirflowFailException}, no retries) instead of writing an empty file.
  \item \textbf{Output:} one JSON file per source with changed content. Large payloads travel between tasks as a temporary file in \texttt{datalake\_bronze/\_tmp/}, not through XCom. Files are written to a \texttt{.part} name and renamed, so the sensor never sees a half-written file.
\end{itemize}

\subsection{\texttt{silver\_processing}}
\begin{itemize}
  \item \textbf{Schedule:} \texttt{15 * * * *} (15 minutes after Bronze), no catch-up, one active run.
  \item \textbf{Trigger:} \texttt{NewBronzeFileSensor}, a \texttt{FileSensor} on \texttt{datalake\_bronze/*.json} that only succeeds when a file is missing from the Silver manifest. A plain FileSensor would succeed on files processed long ago.
  \item \textbf{Sensor settings:} \texttt{poke\_interval}~=~60~s, i.e.\ how often the folder is checked. \texttt{timeout}~=~40~min is how long to wait before giving up, which keeps the sensor shorter than the hourly cycle so two runs never overlap. With \texttt{mode=reschedule} the worker slot is free between pokes. With \texttt{soft\_fail=True}, a run with nothing new ends as skipped, not failed.
  \item \textbf{Task \texttt{process}:} applies Section~\ref{sec:prep} and writes the Parquet files, the quarantine and the state. The state is updated only after all files are written, so a failed run is repeated in full next time.
\end{itemize}

\section{Results evidence}

\shot{airflow_dags.png}{0.95}{Airflow UI: both DAGs active.}
\shot{bronze_grid.png}{0.95}{\texttt{bronze\_ingestion}: successful runs in the grid view.}
\shot{silver_grid.png}{0.95}{\texttt{silver\_processing}: sensor and processing task.}

\subsection{First runs (2026-11-16)}
\begin{tabular}{lrl}
\toprule
\textbf{Bronze file} & \textbf{Records} & \textbf{Size}\\
\midrule
\texttt{gvp\_volcanoes\_20261116\_223828.json} & 165 & 376 KB\\
\texttt{gvp\_eruptions\_20261116\_223856.json} & 11,089 & 11.6 MB\\
\texttt{usgs\_hans\_20261116\_223832.json} & 69 & 40 KB\\
\texttt{usgs\_feed\_20261116\_223830.json} & 1,813 & 1.8 MB\\
\texttt{usgs\_volcano\_updates\_20261116\_223834.json} & 9 & 18 KB\\
\bottomrule
\end{tabular}

\medskip
A second run right after the first wrote only a new feed file: GVP was \emph{not due}, and HANS and the scraped page were \emph{unchanged}. Silver then wrote no new rows, because the new feed file had no new or revised earthquakes near a volcano. Samples of Bronze, Silver and quarantine files are in \texttt{workshop\_2/samples/}, and \texttt{notebooks/02\_silver\_exploration.ipynb} reads them with pandas.

\begin{lstlisting}
>>> pd.read_parquet("datalake_silver/earthquakes_20261116_225940.parquet")[cols].head(3)
         record_id                    event_time   mag mag_type  depth_km  nearest_vnum  distance_km  is_outlier_depth_z                     bronze_file
0  usgs:nc75464572 2026-11-09 22:45:37.890+00:00  1.05       md      0.16        323100        19.17               False  usgs_feed_20261116_224227.json
1  usgs:nc75464587 2026-11-09 23:55:18.600+00:00  1.02       md      4.01        323100        26.88               False  usgs_feed_20261116_225716.json
2  usgs:hv75068917 2026-11-10 01:49:48.400+00:00  2.58       ml     28.73        332010        26.96                True  usgs_feed_20261116_225716.json
\end{lstlisting}

\section{Challenges and decisions}

\begin{itemize}
  \item \textbf{Choosing a scraping target.} The Smithsonian Weekly Report HTML answers HTTP~403 with a bot challenge, so we did not use it. The USGS Volcano Updates page is server-rendered, allowed by \texttt{robots.txt}, US-only, and contains the full notice text.
  \item \textbf{Different notice formats per observatory.} HVO daily updates have no notice number, CVO prints an empty \texttt{VNUM \#}, and AVO uses local time. The parser treats every field as optional, and the record id falls back to the issue time.
  \item \textbf{Global feed versus US scope.} 1,811 events per week, of which 229 fall within 30~km of a US volcano. We apply the Workshop~1 rule (radius and minimum magnitude in \texttt{config/sources.yaml}) and count the rest.
  \item \textbf{Timestamps before 1677} (GVP eruptions) overflow pandas' default nanosecond precision; solved with microsecond timestamps.
  \item \textbf{Repeated quarantine.} The first reruns quarantined the same 129 non-tectonic events every hour. Quarantine is now deduplicated like the data.
  \item \textbf{Failure behaviour (final code review).} A crash between writing the Parquet files and saving the Silver state could duplicate rows on retry; a journal of pending outputs now rolls them back. A Bronze file whose payload breaks a transform is quarantined (\texttt{transform\_error}) instead of blocking every later run. In Bronze, each source runs in its own mapped task group, so an outage of one API does not stop the others, and validation errors fail at once without retries.
  \item \textbf{Windows host.} The container creates a symlink in \texttt{airflow/logs} that Windows cannot read, which broke test collection; pytest is limited to \texttt{tests/}.
\end{itemize}

\begin{thebibliography}{9}
\bibitem{airflow} Apache Software Foundation. \emph{Apache Airflow 2.10 documentation}: TaskFlow API, FileSensor, Docker Compose. \url{https://airflow.apache.org/docs/apache-airflow/2.10.5/}
\bibitem{gvp} Global Volcanism Program (2026). \emph{Volcanoes of the World} (v.~5.4.0; 7 Aug 2026). Smithsonian Institution. \url{https://doi.org/10.5479/si.GVP.VOTW5-2026.5.4}
\bibitem{usgsfeed} U.S. Geological Survey, Earthquake Hazards Program. \emph{Real-time GeoJSON summary feeds}. \url{https://earthquake.usgs.gov/earthquakes/feed/v1.0/geojson.php}
\bibitem{hans} U.S. Geological Survey, Volcano Hazards Program. \emph{Hazard Notification System (HANS) public API}. \url{https://volcanoes.usgs.gov/hans-public/api}
\bibitem{vhp} U.S. Geological Survey, Volcano Hazards Program. \emph{Volcano Updates}. \url{https://www.usgs.gov/programs/VHP/volcano-updates}
\bibitem{iglewicz} Iglewicz, B. and Hoaglin, D. (1993). \emph{How to Detect and Handle Outliers}. ASQC Quality Press.
\bibitem{pandas} The pandas development team. \emph{pandas}. \url{https://pandas.pydata.org}; Apache Arrow, \emph{PyArrow / Parquet}. \url{https://arrow.apache.org/docs/python/}
\bibitem{bs4} Richardson, L. \emph{Beautiful Soup 4 documentation}. \url{https://www.crummy.com/software/BeautifulSoup/bs4/doc/}
\end{thebibliography}

\end{document}
